\scanpage{361}
\uline{normalize} the subgroup $G\subset\widehat{\mathcal T}^{+}(\pi'_0)$
(\uline{N.B.} in fact we have $G\subset\mathcal T^{+}(\pi'_0)$ ; \deleted{\illegible{1}} and $\pi_0$
is reconstructed from $\pi'_0$ and from
this subgroup of $\pi'_0$) --- \deleted{\unsure{I think}}\inserted{we should}
be able to show without too much trouble
that \deleted{\illegible{1}} in (8), $\widetilde{\widehat{\mathcal T}}^{+}(\pi_0)$ and $\widetilde{\widehat{\widehat{\mathcal T}}}(\pi_0)$
are precisely the normalizers of $G$
in $\widehat{\mathcal T}^{+}(\pi'_0)$ and in $\widehat{\widehat{\mathcal T}}(\pi'_0)$ (which
would indeed imply a), besides) ---
but I see no plausible reason
why (b) should be true --- which would mean,
essentially, that there is no more
``arithmetic \unsure{transcendence}'' defined by the
(``very big'' !) $\widehat{\widehat{\mathcal T}}(\pi'_0)$ (modulo $\widehat{\mathcal T}(\pi'_0)^{+}$), than
that defined by the (``much smaller'') group
$\widetilde{\widehat{\widehat{\mathcal T}}}(\pi_0)=\operatorname{Norm}_{\widehat{\widehat{\mathcal T}}(\unsure{\pi_0|\pi'_0})}(G)$ \dots\ Besides these
conditions a) b) are not imperative
for the fundamental conjecture
\unsure{linking} the $\unsure{\pi_{g,\nu}|\Pi_{g,\nu}}$ to $\Gamma_0$ to be \unsure{consistent} ---
\mspage{360}

\scanpage{362}
\mspage{183}
on the other hand we would absolutely need
c) in order at least to define
$A(\pi)\unsure{\longrightarrow|{-}}A(\pi')$ and (if $\Sigma,\Sigma'$ correspond)
$\Gamma_{\Sigma}\unsure{\longrightarrow|{-}}\Gamma_{\Sigma'}$, in the case $\unsure{\pi'_a|\pi'_0}$ \deleted{\unsure{invariant}}\inserted{characteristic} in $\pi$
(and in all cases if in addition we admit
a), which seems \unsure{rather} plausible, and gives
the \uline{injectivities} required. But one
does wonder why an element
of \deleted{\illegible{1}} $\widehat{\widehat{\mathcal T}}(\pi'_0)$ simply because
it \deleted{\unsure{normalizes}}\inserted{is in} the \deleted{$G$ and it \unsure{normalizes}}\inserted{normalizer $R$ of $G$}
\deleted{\unsure{the same}}\inserted{in} $\widehat{\widehat{\mathcal T}}(\pi'_0)$, \inserted{hence also $R$,} and it normalizes
$R\cap\widehat{\mathcal T}^{+}(\pi'_0)$, should also normalize
$\widehat{\mathcal T}^{+}(\pi'_0)$ itself ! It is true that we
made \unsure{draconian} hypotheses at the outset
(\unsure{starting} from a \uline{characteristic} subgroup $\pi'_0$ of $\pi_0$) which
must surely be \unsure{reflected} in special
properties of $G\subset\mathcal T^{+}(\pi'_0)$. Perhaps
in the end the trick of reducing to
characteristic subgroups to examine
\mspage{361}

\scanpage{363}
the situation is \unsure{itself} not so clever.
To see more clearly, we could already
\unsure{try} to understand the case where (without
assuming $\pi'$ characteristic) we assume $\pi'$ of index
2 in $\pi$, hence invariant and $G\simeq\mathbb Z/2$, or
else we take \deleted{\illegible{2}} the kernel of the canonical homomorphism
$\pi\longrightarrow(\pi_{\mathrm{ab}})_2$ ($\simeq(\mathbb Z/2\mathbb Z)^{2g\deleted{+}}$ if $\nu=0$,
$\simeq(\mathbb Z/2\mathbb Z)^{2g\unsure{+|-}\nu-1}$ if $\nu\geq1$), taking for example
$\pi=\widehat\Pi_{0,3}$ --- \unsure{situation} of the
Fermat curve $x^2+y^2+z^2=0$ (octahedral covering
of $\mathbb P^1_{\overline{\mathbb Q}}\setminus\{0,1,\infty\}$ \dots)

But let us provisionally admit the \deleted{results}\inserted{hypotheses}
of injectivity (relatively innocuous), \uline{plus} the
condition c), not innocuous at all ---
whence, if $\Sigma$ and $\Sigma'$ correspond, an injective
homomorphism
\[
(9)\qquad \Gamma_{\Sigma}\hookrightarrow\Gamma_{\Sigma'},
\]
and let us see what we could get from it.
even admitting, provisionally (to
see) that (9) is an isomorphism --- \deleted{\unsure{Which|What}}
that is to say the full force of the main
conjecture $\Gamma_0\simeq\Gamma_{g,\nu}$.
\mspage{362}

\scanpage{364}
\mspage{184}

\unsure{Let|Let} $\pi$ be a profinite group \inserted{with loops} of type $(g,\nu)$, $\pi'$ of type $(g',\nu')$. Let us call a ``correspondence'' between $\pi$ and $\pi'$ a \deleted{group} pair consisting of a profinite $\unsure{G|\pi''}$ \inserted{with loops} \deleted{of type $(g'',\nu'')$}, and of homomorphisms \inserted{with loops} $\unsure{G|\pi''}\xrightarrow{p}\pi$, $\unsure{G|\pi''}\xrightarrow{q}\pi'$ which are therefore \unsure{each|each} \unsure{composed|composed} of a ``plugging holes'' morphism, and of an \deleted{inclusion} isomorphism onto a \unsure{subgroup|subgroup} of finite index.
\unsure{Let|Let} $\unsure{J_p|I''_p}$, $\unsure{J_q|I''_q}\in\unsure{J|I''}=I(\pi'')$ be the \inserted{sub}sets of $\unsure{J|I''}$ which correspond to the ``holes plugged by $p$'' resp. by $q$, we assume \inserted{if necessary} $\unsure{J_p|I''_p}\cap\unsure{J_q|I''_q}=\varnothing$ (otherwise one could factor through the \unsure{same} quotient $\widetilde{\pi}''$ of $\pi''$). In other words if $\unsure{\overline{J}_p|\overline{I}''_p}$, $\unsure{\overline{J}_q|\overline{I}''_q}$ are the complements of $\unsure{J_p|I''_p}$, $\unsure{J_q|I''_q}$ in $\unsure{J|I''}$ (so that we have : $\unsure{\overline{J}_p|\overline{I}''_p}\to I=I(\pi)$, $\unsure{\overline{J}_q|\overline{I}''_q}\to I'=I(\pi')$) we have
\[
\unsure{\overline{J}_p|\overline{I}''_p}\cup\unsure{\overline{J}_q|\overline{I}''_q}=\unsure{J|I''}.
\]

\deleted{\unsure{Using|Let us use} the ``results'' of the two \unsure{paragraphs} \unsure{above}, we then define maps}

We are now going to suppose $\pi$ equipped with a $\Sigma$ \inserted{$\in A(\pi)$}, $\pi'$ equipped with a $\Sigma'\in A(\pi')$, we shall call an ``arithmetic correspondence'' between $\pi$ and $\pi'$ a quadruple $(\unsure{G|\pi''},p,q,\unsure{\Sigma_G|\Sigma''})$, where $(\unsure{G|\pi''},p,q)$ are as above, and $\unsure{\Sigma_G^{a}|\Sigma_G}\in A(G)$, such that we have $\unsure{\Sigma_G|\Sigma_a}=p^*(\Sigma)=q^*(\Sigma')$. \deleted{\unsure{We suppose}} If $P_{\Sigma}$,
\mspage{363}

\scanpage{365}

$P_{\Sigma'}$, $P_{\unsure{\Sigma_G|\Sigma_a}}$ are resp. the torsors under $\Gamma_{\Sigma}$, $\Gamma_{\Sigma'}$, $\Gamma_{\Sigma''}$ that we know, we have two \unsure{diagrams|diagram} of isomorphisms of torsors
\diagram{\begin{tikzcd}[ampersand replacement=\&,column sep=small,row sep=large]
\& P_{\unsure{\Sigma_G|\Sigma_a}} \& \\
P_{\Sigma} \arrow[ur,"p_P^*","\sim"'] \& \& P_{\Sigma'} \arrow[ul,"q_P^*"',"\sim"]
\end{tikzcd}}
\diagram{\begin{tikzcd}[ampersand replacement=\&,column sep=small,row sep=large]
\& \Gamma_{\unsure{\Sigma_G|\Sigma''_G}} \arrow[dr,leftrightarrow,"q_\Gamma^*","\sim"'] \& \\
\Gamma_{\Sigma} \arrow[ur,"\unsure{\widehat{p}_\Gamma|p_\Gamma^*}","\sim"'] \& \& \Gamma_{\Sigma''}
\end{tikzcd}},
whence by composition an \unsure{isomorphism|homomorphism} of
\deleted{$\unsure{P_{\Sigma}\to}$}
\[
(P_{\Sigma},\Gamma_{\Sigma})\xrightarrow{\sim}(P_{\Sigma'},\Gamma_{\Sigma'})
\]

We call ``arithmetically outer isomorphisms'' \inserted{of $(\pi,\Sigma)$ with $(\pi',\Sigma')$,} any isomorphism of torsors that we can obtain in this way. Two correspondences are said to be \inserted{\unsure{equivalent} \deleted{\unsure{from the point of}} from the point of view} \deleted{\unsure{arithmetic|arithmetic} \unsure{outer|outer}} arithmetically outer, if they define the \unsure{same|=} arithmetically outer isomorphism. \unsure{The} composition is defined by composition of the actions on $(P_{\Sigma},\Gamma_{\Sigma})$ --- we see that \unsure{this} comes from a correspondence. We \unsure{find} therefore a groupoid, which we check without difficulty is connected. \inserted{I say that} The automorphisms of $(\pi,\Sigma)$ ``are'' the automorphisms of $(P_{\Sigma},\Gamma_{\Sigma})$ defined by elements of $\Gamma_{\Sigma}$. \deleted{This is \unsure{obvious} or \unsure{not}} This is \unsure{easy}\ldots

\margin{We shall call it the \unsure{groupoid} of ``outer arithmetic curves''.}
\mspage{364}

\scanpage{366}
\mspage{185}

In this groupoid, there is a transitive system of isomorphisms between the \deleted{\illegible{1}} $(\widehat{\pi}_0,\overline{\Sigma}_{\pi_0})$ (where $\pi_0$ is a discrete group with loops) more generally between the $(\pi,\Sigma_\alpha)$, where $\alpha$ \inserted{$\in P(\pi)$} is a prediscretification of $\pi$ --- let us therefore define an \inserted{arithmetically} outer isomorphism of $\Pi_{g,\nu}$
$\pi$ \ldots\ \unsure{The|Their} group of \unsure{these|its} automorphisms is denoted $\Gamma_0$. The \deleted{category} groupoid of \inserted{virtual} outer arithmetic curves is therefore identified with the category of torsors under $\Gamma_0$,
A prediscretification (we could also call it an arithmetically outer rigidification) is \unsure{therefore} nothing other than an arithmetically outer isomorphism between this reference element, and $\pi$.
\mspage{365}

\scanpage{367}
\mspage{186}

\uline{29} Critique of the preceding approach.

The approach \unsure{of the|of the} preceding \S\ seems \unsure{frankly|in the end} very brutal. I even have doubts whether \unsure{the|the} \deleted{hypothesis} conjecture \unsure{in|on} the properties of the ``plugging holes'' functor is true as it stands. It is true that for a profinite group \deleted{\illegible{1}} \inserted{\unsure{with}} \unsure{loops} $\pi$, if we fix $i\in I(\pi)$, we have a canonical homomorphism
\[
\tag{1}
\widehat{\widehat{\mathcal T}}(\pi,i)
=\underset{\substack{\text{stabilizer}\\\text{of }i\text{ in }\widehat{\widehat{\mathcal T}}(\pi)}}{\text{\uline{$\widehat{\widehat{\mathcal T}}(\pi)_i$}}}
\longrightarrow \widehat{\widehat{\mathcal G}}(\pi'),
\]
but it is problematic whether this is an isomorphism --- and even whether it is injective, or whether it is surjective. But, \unsure{choosing} a discretification $\pi_0\subset\pi$, whence ditto $\pi'_0$ for $\pi'$, the preceding homomorphism does indeed induce an isomorphism
\[
\tag{1}
\mathcal T(\pi_0)_i\simeq\mathcal G(\pi'_0)
\]
whence
\[
\tag{3}
\widehat{\mathcal T}(\pi_0)\xrightarrow{\sim}\widehat{\mathcal G}(\pi'_0)
\]
and consequently, the exact sequence $\deleted{\unsure{\widehat{\pi}_0}}\widehat{\pi}'_0\longrightarrow\widehat{\mathcal G}(\pi'_0)\longrightarrow\widehat{\mathcal T}(\pi'_0)\longrightarrow 1$ gives
\[
\tag{4}
1\longrightarrow\widehat{\pi}'_0\longrightarrow\widehat{\mathcal T}(\pi_0)\longrightarrow\widehat{\mathcal T}(\pi'_0)\longrightarrow 1
\]
and consequently we find an injective homomorphism
\[
\begin{gathered}\widehat{\pi}'_0\\[-0.5ex]\Vert\\[-0.5ex]\pi'\end{gathered}
\unsure{\hookrightarrow|\leftrightarrow}
\widehat{\mathcal T}(\pi_0)_i
\hookrightarrow
\begin{gathered}\widehat{\widehat{\mathcal T}}(\pi_0)_i\\[-0.5ex]\Vert\\[-0.5ex]\unsure{\widehat{\widehat{\mathcal T}}}(\pi)\end{gathered}
\]
hence a homomorphism
\[
\tag{5}
i_{\pi_0}:\quad\pi'\longrightarrow\widehat{\widehat{\mathcal T}}(\pi)_i
\]
for which \unsure{the|the} composite with (1) is the canonical injection $\pi'\hookrightarrow\widehat{\widehat{\mathcal G}}(\pi')$. Replacing the discretification \inserted{$\pi_0$} by another, $\pi_1$, we find a priori another homomorphism $i_{\pi_1}:\pi'\longrightarrow\widehat{\widehat{\mathcal T}}(\pi)_i$. \unsure{We may assume} $\pi_1=u(\pi_0)$, $u\in\widehat{\widehat{\mathcal T}}(\pi)_i$;

\mspage{366}

\scanpage{368}

and by transport of structure we find
\[
i_{\pi_1}=i_{u(\pi_0)}=\operatorname{int}(u)\circ i_{\pi_0}\circ u_{\pi'}^{-1}
\]
where $u_{\pi'}$ is the automorphism of $\pi'$ defined by $u$ via (1). To say that $i_{\pi_0}$ is independent of the choice of the discretification $\pi_0$, also amounts to saying that its image in $\widehat{\widehat{\mathcal T}}(\pi)_i$ is an invariant subgroup --- and then the action of $\widehat{\widehat{\mathcal T}}(\pi)_i$ on \deleted{\unsure{$\pi$}} the invariant subgroup $i_{\pi_0}(\pi)=i(\pi)$ \deleted{\illegible{1}} \unsure{via} inner \unsure{automorphisms}, is none other than that defined by (1). \uline{In} this case, we find by passage to the quotient a \unsure{homomorphism}
\[
\tag{6}
\widehat{\widehat{\mathcal T}}(\pi)_i/\pi'\longrightarrow\widehat{\widehat{\mathcal T}}(\pi')
\]
whose injectivity resp. surjectivity would be equivalent to that of (1). But it is not at all obvious that it is \unsure{always} so.

If it were not so, it \deleted{would be} would be necessary to look at the subgroup of $\widehat{\widehat{\mathcal T}}(\pi)_i$ \inserted{formed of the $\gamma\in\unsure{\widehat{\widehat{\mathcal T}}}(\pi)_i$ which} \deleted{which} normalize the \deleted{\unsure{injecting|injected}} subgroup $i_{\pi_0}(\pi')$, and such that the action induced on $i_{\pi_0}(\pi')$ corresponds to that given by the action (1) of $\widehat{\widehat{\mathcal T}}(\pi)$ on $\pi'$. This subgroup \inserted{$H_{\pi_0,\unsure{i}}$ (which contains $\widehat{\mathcal T}(\pi_0)_i$) \deleted{\illegible{1}}} therefore depends a priori on the chosen discretification. Replacing $\pi_0$ by $u(\pi_0)$ (where \deleted{\unsure{$u$}} $u\in\widehat{\widehat{\mathcal T}}(\unsure{\pi_0|\pi})_i$) replaces it by $\operatorname{int}(u).H$ --- hence $H$

\mspage{367}

\scanpage{369}
\mspage{187}

at least does not change, if we vary $\pi_0$ within \unsure{their|a} class of prediscretification.

We can do more symmetric things, by considering \uline{for every} $i\in I$ the \unsure{corresponding} quotient $\pi'_i$ of $\pi$, whence, for a chosen discretification $\pi_0$, homomorphisms
\[
\unsure{i_{\pi_0,i}}:\ \pi'_i\longrightarrow
\widehat{\widehat{\mathcal T}}(\pi)_i\subset
\widehat{\widehat{\mathcal T}}(\pi)
\]
and we shall define $H_{\pi_0}\subset\widehat{\widehat{\mathcal T}}(\pi)$ as the subgroup of the $\gamma\in\widehat{\widehat{\mathcal T}}(\pi)$ which ``permute the \deleted{\illegible{1}}\inserted{$i_{\pi_0,i}$} among themselves'' in an evident sense. Thus we have
\[
\tag{7}
H_{\pi_0}\cap\widehat{\widehat{\mathcal T}}(\pi)_i=H_{\pi_0,i}
\]
and
\[
\tag{8}
H\supset\widehat{\mathcal T}(\pi_0)
\]
The point I have in view, is that the subgroup of $\widehat{\widehat{\mathcal T}}(\pi_0)$ image of $\pi_1(M_{g,\nu})$, must not only normalize $\widehat{\mathcal T}(\pi_0)^+$, but moreover be contained in this $H$. This is a condition I have \unsure{encountered} by the back door, \deleted{\unsure{under}}\inserted{thanks to} the hypothesis (\unsure{may} be very hasty) that (1) is an isomorphism, which implied (tacitly) that we had $H=\widehat{\widehat{\mathcal T}}(\pi_0)$ in its entirety. It is possible that $\widehat{\widehat{\mathcal T}}(\pi_0)$ is a group so \unsure{enormous|unreasonable} and pathological, that there will never be any question of saying reasonable \inserted{(and true)} things about \deleted{\unsure{the image of}} this whole group (such as the bijectivity of (1), for example

\mspage{368}

\scanpage{370}

and that we are obliged to work with smaller subgroups, which \unsure{remain} close to the discrete (with all the same some additional ``arithmetic'' aspects, due to the $\Gamma_0=\mathrm{Gal}(\unsure{\overline{\mathbb Q}_0}/\mathbb Q)$ !). In fact, there is \inserted{(for $I=I(\pi)\ne\varnothing$ i.e.\ $\nu\ne0$)} in the group $\widehat{\mathcal T}(\pi_0)$ a \inserted{\unsure{simplicial}} structure of successive extensions, which is going to be respected by the outer action of the Galois group, and which we should take into account. It is part of the ``structure \unsure{at|:} $\infty$'' in the $\pi_1$ of the \deleted{\unsure{spaces}} modular multiplicities $M_{g,\nu}$, which even for $\nu=0$ is probably nontrivial, and it may be necessary to take it into account, to manage to put our finger on $\Gamma_{\mathbb Q}$.

Without trying to give from the outset an a priori description of $\Gamma_{\mathbb Q}$ ``in the $\widehat{\widehat{\mathcal T}}_{g,\nu}$'', we are going to proceed in a more inductive way, starting from the presence of $\Gamma_{\mathbb Q}$ (for \unsure{arithmetico-geometric} reasons), and trying to bring out properties of this presence, perhaps strong enough to end up giving a purely algebraic characterization.

\mspage{369}
